% !TEX root = main.tex
\section{Strategie}
\label{set:strategie}

   \subsection{Inverse Volatility}
\label{sec:inverse_volatility}

L'idea fondamentale di questa strategia consiste nell'assegnare il capitale in modo \textbf{inversamente proporzionale} alla volatilità di ciascun asset, penalizzando i titoli più nervosi e premiando quelli più stabili.

Il calcolo dell'allocazione si articola in tre passaggi principali:
\begin{enumerate}
    \item \textbf{Misura del rischio}: Per ogni asset $i$, si calcola la volatilità storica $\sigma_i$ applicando la deviazione standard campionaria sui rendimenti.
    \item \textbf{Inversione}: Si calcola il reciproco della volatilità $\frac{1}{\sigma_i}$. Se un titolo presenta un'elevata volatilità, il suo inverso risulterà contenuto; viceversa, per un titolo stabile il reciproco assume un valore più elevato.
    \item \textbf{Normalizzazione}: Si divide l'inverso della volatilità del singolo asset per la somma degli inversi di tutti gli asset del portafoglio, garantendo che la somma dei pesi $w_i$ sia pari a $1$ ($100\%$ del capitale allocato).
\end{enumerate}

La formula matematica per il calcolo del peso $w_i$ dell'asset $i$-esimo è definita come:

\begin{equation}
    w_i = \frac{\frac{1}{\sigma_i}}{\sum_{j=1}^{N} \frac{1}{\sigma_j}}
\end{equation}

dove:
\begin{itemize}
    \item $w_i$ è il peso percentuale assegnato all'asset $i$-esimo nel portafoglio;
    \item $\sigma_i$ è la deviazione standard campionaria dei rendimenti dell'asset $i$;
    \item $N$ è il numero totale di asset presenti nel portafoglio;
    \item $\sum_{j=1}^{N} \frac{1}{\sigma_j}$ rappresenta la somma degli inversi delle volatilità di tutti gli asset (fattore di normalizzazione).
\end{itemize}

Nel codice TypeScript riportato, vengono inseriti dei controlli difensivi: se la deviazione standard di un asset è pari a zero ($\sigma_i = 0$), il suo valore reciproco viene impostato a $0$ per evitare la divisione per zero. Allo stesso modo, se la somma totale degli inversi risulta nulla, la funzione restituisce un vettore di pesi azzerati.

\begin{lstlisting}[language=TypeScript, caption={Implementazione dell'algoritmo Inverse Volatility.}, label={lst:inverse_volatility}]
export function inverseVolatility(returnsByAsset: number[][]): number[] {
  // 1. Calcolo del reciproco della volatilita per ogni asset
  const inv = returnsByAsset.map((r) => {
    const s = sampleStd(r);
    return s === 0 ? 0 : 1 / s; // Guard divisione per zero
  });

  // 2. Somma totale dei reciproci
  const tot = inv.reduce((a, b) => a + b, 0);
  if (tot === 0) return inv.map(() => 0);

  // 3. Normalizzazione a somma 1 (100% del capitale)
  return inv.map((x) => x / tot);
}
\end{lstlisting}

    \subsection{Kelly Criterion}
\label{sec:kelly_criterion}

Il criterio di Kelly calcola l'esposizione ottimale di capitale per massimizzare il tasso di crescita logaritmico del portafoglio nel tempo. La formula classica in forma matriciale è definita come:

\begin{equation}
    \mathbf{f}^* = \mathbf{\Sigma}^{-1} \boldsymbol{\mu}
\end{equation}

\begin{itemize}
    \item $\mathbf{f}^*$ è il vettore dei pesi ideali delle posizioni;
    \item $\mathbf{\Sigma}^{-1}$ è l'inversa della matrice di covarianza dei rendimenti;
    \item $\boldsymbol{\mu}$ è il vettore dei rendimenti attesi (\textit{expected returns}).
\end{itemize}

A questa formulazione teorica vengono applicate quattro modifiche pratiche. La formula pura, infatti, risulta estremamente sensibile agli errori di stima degli input (\textit{garbage in, garbage out}) e tende a generare una leva finanziaria eccessivamente aggressiva ed instabile.

Le quattro modifiche apportate sono le seguenti:

\begin{enumerate}
    \item \textbf{Regolarizzazione della Matrice (\textit{Covariance Shrinkage})}: Se due o più titoli presentano un comportamento fortemente correlato o se il numero di campioni storici è ridotto, la matrice di covarianza diventa singolare (o quasi singolare). L'inversione di una tale matrice in JavaScript comporta gravi errori di precisione numerica o il blocco dell'applicazione. Per garantire la stabilità computazionale, viene applicata una regolarizzazione sui parametri fuori dalla diagonale principale:
    \begin{equation}
        \mathbf{\Sigma}_{\text{shrunk}}(i, j) = 
        \begin{cases} 
            \mathbf{\Sigma}(i, j) & \text{se } i = j \quad \text{(diagonale principale: varianze inalterate)} \\ 
            (1 - \delta) \cdot \mathbf{\Sigma}(i, j) & \text{se } i \neq j \quad \text{(fuori diagonale: covarianze ridotte del } 15\% \text{)}
        \end{cases}
    \end{equation}
    In questo modo, i titoli appaiono leggermente meno correlati tra loro, garantendo un'inversione matriciale numericamente stabile.

    \item \textbf{Scalamento Half-Kelly}: Dopo aver calcolato il vettore $\mathbf{f}^*$, si applica uno scalamento dimezzandone i valori:
    \begin{equation}
        \mathbf{f}_{\text{half}} = \frac{1}{2} \mathbf{f}^*
    \end{equation}
    Lo scopo è mitigare i \textit{drawdown} estremi causati da una leva eccessiva. Dimezzare le posizioni riduce la varianza complessiva del portafoglio del $75\%$, sacrificando solo una quota minima del tasso di crescita teorico.

    \item \textbf{Vincolo Long-Only}: La formula originale restituisce sia valori positivi che negativi (dove un valore positivo come $+0.4$ indica l'acquisto del titolo col $40\%$ del capitale, mentre un valore negativo come $-0.25$ suggerisce la vendita allo scoperto per il $25\%$). Poiché lo \textit{short selling} richiede garanzie sui margini e autorizzazioni complesse, si è scelto di eliminare le posizioni corte applicando una funzione di taglio:
    \begin{equation}
        f_i^{\text{long}} = \max\left(0, f_{\text{half}, i}\right)
    \end{equation}

    \item \textbf{Cap della Leva Massima (\textit{Leverage Ceiling})}: Il criterio di Kelly tende ad adottare una strategia di leva molto aggressiva. Si è stabilito quindi un limite massimo alla leva finanziaria $L_{\max} = 2$. Se l'esposizione lorda complessiva $L = \sum f_i^{\text{long}}$ supera tale soglia, i pesi vengono riallocati proporzionalmente:
    \begin{equation}
        \mathbf{f}_{\text{finale}} = 
        \begin{cases} 
            \mathbf{f}^{\text{long}} & \text{se } L \le L_{\max} \\ 
            \mathbf{f}^{\text{long}} \cdot \frac{L_{\max}}{L} & \text{se } L > L_{\max} 
        \end{cases}
    \end{equation}
\end{enumerate}

Nel Listing viene riportata l'implementazione in TypeScript della funzione descritta.

\begin{lstlisting}[language=TypeScript, caption={Implementazione dell'algoritmo di Kelly Multi-Asset.}, label={lst:kelly_criterion}]
export function kelly(
  mu: number[],
  cov: number[][],
  opts: { shrink?: number; halfKelly?: boolean; maxLeverage?: number } = {},
): number[] {
  const { shrink = 0.15, halfKelly = true, maxLeverage = 2 } = opts;

  // 1. Regolarizzazione della matrice (Shrinkage)
  const shrunk = cov.map((row, i) =>
    row.map((v, j) => (i === j ? v : (1 - shrink) * v)),
  );
  const covInv = invertMatrix(shrunk);

  // 2. Calcolo f* = Sigma^-1 * mu
  let f = matVec(covInv, mu); 

  // 3. Vincolo Long-Only
  f = f.map((x) => Math.max(0, x)); 

  // 4. Cap della Leva Massima
  const gross = f.reduce((a, b) => a + b, 0);
  if (gross > maxLeverage) f = f.map((x) => (x * maxLeverage) / gross); 

  // 5. Half-Kelly
  if (halfKelly) f = f.map((x) => x / 2); 

  return f; // Esposizioni (il capitale non allocato rimane in Cash)
}
\end{lstlisting}

\subsection{Target Volatility}
\label{sec:target_volatility}

L'obiettivo della strategia di \textit{Target Volatility} (o \textit{Volatility Targeting}) è controllare il rischio complessivo del portafoglio, mantenendo la sua volatilità stimata entro una soglia predefinita ($\sigma_{\text{target}}$). Se la volatilità del portafoglio di partenza supera tale limite, l'algoritmo riduce proporzionalmente l'esposizione su tutti i titoli, allocando la quota rimanente di capitale in liquidità (\textit{Cash}).

L'implementazione segue tre passaggi principali:

\begin{enumerate}
    \item \textbf{Volatilità del portafoglio (giornaliera e annualizzata)}: Si calcola prima la deviazione standard giornaliera mediante la forma quadratica della matrice di covarianza $\mathbf{\Sigma}$ e del vettore dei pesi base $\mathbf{w}_{\text{base}}$, per poi annualizzarla moltiplicando per la radice dei giorni lavorativi ($\sqrt{252}$):
    \begin{equation}
        \sigma_{\text{port, giornaliera}} = \sqrt{\mathbf{w}_{\text{base}}^T \mathbf{\Sigma} \mathbf{w}_{\text{base}}}
    \end{equation}
    \begin{equation}
        \sigma_{\text{port, annua}} = \sigma_{\text{port, giornaliera}} \times \sqrt{252}
    \end{equation}

    \item \textbf{Fattore di scalamento}: Viene calcolato il coefficiente di riduzione $\text{scale}$:
    \begin{equation}
        \text{scale} = \min\left(1, \; \frac{\sigma_{\text{target}}}{\sigma_{\text{port, annua}}}\right)
    \end{equation}
    L'utilizzo dell'operatore $\min(1, \dots)$ impedisce l'impiego della leva finanziaria, evitando che l'algoritmo aumenti l'esposizione complessiva oltre il $100\%$ del capitale nei periodi di bassa volatilità. Inoltre, il rapporto $\frac{\sigma_{\text{target}}}{\sigma_{\text{port, annua}}}$ esegue un de-risking automatico nelle fasi di elevata turbolenza di mercato: quando la volatilità annua registra picchi improvvisi, il fattore $\text{scale}$ si riduce automaticamente, vendendo quote dei titoli per rifugiarsi nella liquidità e proteggere il portafoglio da ampi \textit{drawdown}.

    \item \textbf{Pesi scalati}: Ogni peso dell'allocazione base viene ridimensionato applicando il fattore di scalamento:
    \begin{equation}
        \mathbf{w}_{\text{finale}} = \mathbf{w}_{\text{base}} \times \text{scale}
    \end{equation}
\end{enumerate}

La funzione gestisce anche il caso limite in cui la volatilità annualizzata sia nulla ($\sigma_{\text{port, annua}} = 0$), impostando il fattore di scala a $1$ per evitare divisioni per zero.

\begin{lstlisting}[language=TypeScript, caption={Implementazione dell'algoritmo Target Volatility.}, label={lst:target_volatility}]
export function targetVolatility(
  wBase: number[],
  cov: number[][],
  sigmaTargetAnnual = 0.097, // 2.8% mensile ~ 9.7% annuo
): number[] {
  // 1. Calcolo della volatilita giornaliera ed annualizzata del portafoglio
  const sigmaDaily = Math.sqrt(quadForm(wBase, cov));
  const sigmaAnnual = sigmaDaily * Math.sqrt(252);

  // 2. Calcolo del fattore di scalamento (cap a 1, no leva)
  const scale =
    sigmaAnnual === 0 ? 1 : Math.min(1, sigmaTargetAnnual / sigmaAnnual);

  // 3. Applicazione del fattore ai pesi (il capitale rimanente va in Cash)
  return wBase.map((w) => w * scale);
}
\end{lstlisting}